\documentclass[11pt]{article}
\usepackage[T1]{fontenc}
\usepackage[utf8]{inputenc}
\usepackage{lmodern,microtype}
\usepackage[margin=0.86in,headheight=15pt]{geometry}
\usepackage{amsmath,amssymb,amsthm,mathtools}
\usepackage{booktabs,array,longtable,enumitem,xcolor}
\usepackage[numbers,sort&compress]{natbib}
\usepackage{fancyhdr}
\usepackage[colorlinks=true,linkcolor=blue!45!black,citecolor=blue!45!black,urlcolor=blue!45!black]{hyperref}
\usepackage{listings}
\lstset{basicstyle=\ttfamily\small,columns=fullflexible,breaklines=true,frame=single,rulecolor=\color{black!20}}
\pagestyle{fancy}
\fancyhf{}
\fancyhead[L]{\small Exact temporal walk profiles}
\fancyhead[R]{\small Review manuscript / 12 November 2026}
\fancyfoot[C]{\thepage}
\setlength{\parskip}{0.35em}
\setlength{\parindent}{0pt}
\setlist{nosep,leftmargin=*}
\newtheorem{theorem}{Theorem}[section]
\newtheorem{lemma}[theorem]{Lemma}
\newtheorem{corollary}[theorem]{Corollary}
\newtheorem{proposition}[theorem]{Proposition}
\theoremstyle{remark}
\newtheorem{remark}[theorem]{Remark}
\newcommand{\cost}{\operatorname{cost}}
\newcommand{\poly}{\operatorname{poly}}
\newcommand{\N}{\mathbb{N}}
\newcommand{\infv}{\infty}
\newcommand{\Ppart}{P_{\le R}}
\title{Minimum-Cost Temporal Walk Profiles\\[0.15em]
on Depth-Three In-Arborescences\\[0.4em]
\large Exact optimization parameterized by preservation loss}
\author{Goutam Dyapa}
\date{Working manuscript --- 12 November 2026}
\hypersetup{pdfauthor={Goutam Dyapa},pdftitle={Minimum-Cost Temporal Walk Profiles on Depth-Three In-Arborescences},pdfsubject={Exact optimization parameterized by preservation loss}}

\begin{document}
\maketitle
\begin{abstract}
This paper studies minimum-cost deletion of explicitly listed temporal events subject to
preserving a prescribed number of two-event walks ending at a designated root.
For temporal in-arborescences of depth at most three, the algorithm computes the minimum
deletion cost for every attainable exact pair of surviving two- and three-event
walk counts. Events have nonnegative binary-encoded costs and may be required
to survive or to be deleted. If $R$ is the allowed loss of two-event walks after mandatory deletions,
the algorithm runs in $2^{O(\sqrt R)}\poly(L)$ time and polynomial space,
where $L$ is the complete encoded input length. The unit-cost formulation
admits a polynomial-size generalized deletion propagation encoding, and
binary deletion costs admit an explicit weighted ILP extension. The result
concerns the specialized algorithm and its dependence on preservation loss. The argument
charges deleted final events to preservation loss, groups them by
continuation-equivalent neighborhoods, and enumerates a partition-bounded
family of branch configurations. Chronological dynamic programs recover exact
profiles and original-event witnesses. For contrast, the corresponding
unit-cost quota decision problem is NP-complete and $\mathrm{W}[1]$-hard in
the deletion budget on the same graph class. The profiles support arbitrary
efficiently evaluated predicates of the two counts and deletion-invariant
data; an elementary hardness example delineates an additional dynamic-feature
boundary. An exhaustive check on small instances compares the reference
implementation with direct enumeration. A study of synthetic
transaction histories illustrates why a three-walk upper bound and the
corresponding exact target can require different answers.
\end{abstract}

\section{Introduction}\label{sec:intro}

Deleting a temporal event can destroy several time-respecting interactions
at once. A decision may therefore have two competing effects: it can suppress
longer interactions while also removing shorter ones that are required to
survive. The model captures this tradeoff on a directed in-arborescence of depth at most three, with all walks
ending at its root. Parallel contacts are separate events, and successive
timestamps must increase strictly. The algorithm chooses a shared set of
event deletions, preserving a lower quota on two-event walks and recording
exactly which two- and three-event walk counts can remain.

Temporal reachability modification and path counting already have substantial
algorithmic foundations. Molter, Renken, and
Zschoche~\citep[Theorems 14--15]{molter2024} give polynomial forest
algorithms for multiple-source temporal reachability minimization by deletion
or delay, and Enright, Meeks, and Molter give a polynomial forest algorithm
for counting temporal paths~\citep{enright2025counting}.
Generalized deletion propagation (GDP) also provides a framework for
simultaneous deletion and preservation requirements over query
outputs~\citep{makhija2025gdp}. Appendix~\ref{app:gdp} gives a polynomial-size
encoding of the present unit-cost formulation in GDP and an explicit
weighted ILP extension for binary deletion costs.

The contribution here is the specialized temporal algorithm and its
preservation-loss bound. Its sufficient state preserves the multiplicities
needed for exact two- and three-event counts while optimizing one shared
set of original-event deletions. The complete profile can also be obtained
through polynomially many exact target optimizations; producing that table
alone is therefore not the algorithmic distinction. Section~\ref{sec:prior}
separates these representation issues from the parameter dependence required
to derive Theorem~\ref{thm:main}.

\paragraph{The preservation parameter.}
After applying mandatory deletions, let $N_2$ be the baseline number of
two-event walks and let $T_0$ be the nonnegative preservation quota. The allowed
additional loss is $R=N_2-T_0$. The algorithm computes every admissible exact
count pair, its minimum nonnegative binary-encoded deletion cost, and an
attaining witness in $2^{O(\sqrt R)}\poly(L)$ time and polynomial space
(Theorem~\ref{thm:main}). The polynomial exponent is absolute. Neither a
monetary value nor the number of branches enters the exponential factor.
The input is an explicit event list; binary multiplicities standing for
unlisted events are outside the model.

\paragraph{Where the bound comes from.}
A final event is a contact incident to the root. Its deletion destroys a
number of baseline two-event walks equal to the number of earlier middle
events. These charges are disjoint across final events. Within one branch,
equal charges identify equal prefix neighborhoods and therefore equal effects
under every continuation on the other event layers. Nonnegative costs allow a
canonical cheapest choice within each class. The possible charged
multiplicities are bounded by the number of integer partitions of numbers at
most $R$. The algorithm enumerates these configurations separately for each branch.
For a fixed configuration, a chronological state records the retained
first-event prefix count, exact two-event loss, and exact three-event count.
Exact convolutions then assemble child and branch profiles. This order of
operations keeps one partition factor rather than a product over branches.

\paragraph{Why exact profiles matter.}
A scalar optimum for a monotone quota can discard intermediate count values
that are essential for a nonmonotone target. Moreover, two unfinished choices
can have the same current cost and walk counts yet admit different optimal
continuations. Proposition~\ref{prop:unfinished} gives a ten-event example
where keeping either one of two scalar-equal contexts loses a profile entry
or its minimum cost. The state keeps the information needed until it is safe
to merge completed tables. Retaining exact coordinates, minimizing values,
and min-plus convolution are standard techniques; the structural contribution
is their sufficient temporal state and the preservation-loss bound proved
here.

\paragraph{Complexity boundaries and supported decisions.}
The unit-cost quota problem is $\mathrm{W}[1]$-hard in the deletion budget
even when every event is optional (Theorem~\ref{hard:budget}); its classical
decision version is NP-complete (Corollary~\ref{full:np-complete}). This is a
contrast between parameters on the same graph class, not a lower bound in
$R$. Once a full profile is available, any efficiently evaluated acceptance
predicate of the two counts and invariant auxiliary data can be optimized by
a scan. This includes nonmonotone count-based classifiers. An arbitrary
additional dynamic amount feature already permits NP-completeness at $R=0$
(Proposition~\ref{prop:amount}), so exactness does not extend automatically to
a richer feature pipeline.

The remaining sections give the full structural proof, the companion hardness
result, concrete profile examples, related-work comparisons, and validation.
The complete budget-hardness construction and a counting-model normalization
are included in the appendices. The implementation evidence concerns
correctness on a finite corpus; no practical speed or application-effectiveness
claim is needed for the formal result.

\section{Model and main result}\label{sec:model}

\subsection{Events, counts, and normalization}

Let $E$ be a set of $M$ explicitly listed events, each with a unique original
identifier, a source vertex, a target vertex, an integer timestamp $t(e)$, and
a nonnegative integer deletion cost $c_e$. Costs and timestamps are encoded
in binary for the complexity statement. Parallel events are distinct decisions,
including events with identical endpoints and timestamps. After parallel arcs
are collapsed, the original graph is an in-arborescence directed toward a
designated root $z$, with depth at most three. Every mentioned vertex has a path
to $z$. The empty history is allowed. Depth and roles refer to this original
graph, even if subsequent mandatory deletions disconnect it.

For $S\subseteq E$ and $h\in\{2,3\}$, define $W_h(S)$ as the number of ordered
$h$-tuples of events of $S$ that form a directed walk \emph{ending at $z$}, with
strictly increasing timestamps. Walks ending elsewhere are excluded. In particular, a path of
three increasing events has $W_2=1$ and $W_3=1$, although it has two length-two
subwalks if endpoints are unrestricted.

Let $T$ be an integer preservation quota, and let $F$ and $K$ be disjoint
forced-deleted and forced-retained sets. All other
events are optional. Set
\[
 E_0=E\setminus F,\qquad O=E_0\setminus K,\qquad
 N_2=W_2(E_0),\quad N_3=W_3(E_0),\quad T_0=\max\{T,0\}.
\]
If $T_0>N_2$, no deletion can meet the lower quota and the instance is rejected.
Otherwise let
\begin{equation}
 R=N_2-T_0,\qquad
 \lambda(D_0)=N_2-W_2(E_0\setminus D_0),\quad D_0\subseteq O.
 \label{eq:parameter}
\end{equation}
Then $0\le R\le M^2$, and the quota is exactly $\lambda(D_0)\le R$.
Forced cost $c(F)=\sum_{e\in F}c_e$ is reserved once. It is never charged again
inside a dynamic-programming transition. The input length $L$ includes the
complete explicit encodings of identifiers, graph, timestamps, statuses, costs,
and quota.

\subsection{Exact minimum-cost profiles}

For $0\le\ell\le R$ and $0\le w\le N_3$, define
\begin{equation}
 C(\ell,w)=\min\left\{
 \sum_{e\in D_0}c_e:
 D_0\subseteq O,\quad
 W_2(E_0\setminus D_0)=N_2-\ell,\quad
 W_3(E_0\setminus D_0)=w
 \right\},
 \label{eq:profile}
\end{equation}
with value $\infv$ if the set is empty. A finite entry must include an attaining
set of original identifiers. Its total cost is $c(F)+C(\ell,w)$ and its full
deletion witness is $F\cup D_0$.

\begin{theorem}[Weighted exact attainable profiles]\label{thm:main}
For the model above, all finite entries of~\eqref{eq:profile}, their minimum
costs, and attaining original-event witnesses can be computed in
\[
 2^{O(\sqrt R)}(L+1)^{O(1)}
\]
time and polynomial space. The polynomial exponent does not depend on $R$,
the number of branches, or the numerical magnitudes of the deletion costs.
Strict timestamp ties, parallel events, and both forced statuses are supported.
The output contains at most $(R+1)(N_3+1)$ entries.
\end{theorem}

The output is complete over the \emph{count pairs}, with one minimum-cost
witness per pair. It does not list every attainable monetary cost or every
witness. With unit costs, selecting the cheapest entry satisfying $w\le U$
solves the corresponding upper/lower quota problem. A total monetary budget can be imposed
by filtering entries. An additional, independent cardinality cap in a weighted
instance is not represented by this table; it would require additional state.

\section{Final-event structure and bounded enumeration}\label{sec:structure}

For each depth-one vertex $u$, let $C_u$ be the events $u\to z$ in $E_0$.
For each depth-two child $v$ of $u$, let $B_v$ be the events $v\to u$ in $E_0$,
and let $A_v$ contain all events in $E_0$ from depth-three children into $v$. Put
$B_u=\bigcup_{v\to u}B_v$. These event domains are disjoint. Original children
remain present in the decomposition even if their surviving event lists are
empty. For a retained $b\in B_v$, let
\[
 p_b=|\{a\in A_v\setminus D_0:t(a)<t(b)\}|,\qquad
 s_b=|\{c\in C_u\setminus D_0:t(b)<t(c)\}|.
\]
The unique middle event identifies each counted walk, so
\begin{equation}
 W_2(E_0\setminus D_0)=\sum_u\sum_{v\to u}\sum_{b\in B_v\setminus D_0}s_b,
 \qquad
 W_3(E_0\setminus D_0)=\sum_u\sum_{v\to u}\sum_{b\in B_v\setminus D_0}p_bs_b.
 \label{eq:identities}
\end{equation}
Deleting an $A$ event cannot affect $W_2$ under the root-ending definition.

\begin{lemma}[Final-event charge and equal-degree equivalence]\label{lem:c}
For $c\in C_u$, write
\[
 N(c)=\{b\in B_u:t(b)<t(c)\},\qquad d(c)=|N(c)|.
\]
Every optional deletion set satisfies
\begin{equation}
 \sum_u\sum_{c\in D_0\cap C_u}d(c)\le\lambda(D_0).
 \label{eq:ccharge}
\end{equation}
Within one branch, $C$ events of equal degree have identical neighborhoods.
For any fixed decisions on all $A$ and $B$ events, the two walk counts depend
only on the number of deleted optional $C$ events of each degree, not on their
identities.
\end{lemma}

\begin{proof}
Deleting $c$ destroys the $d(c)$ baseline pairs $(b,c)$. Pairs charged to
different final events are disjoint, including when some middle events are
also deleted. They are a subset of all destroyed baseline pairs, proving the
inequality. Sort $B_u$ by timestamp and break ties in any fixed way. Each
$N(c)$ is a prefix of that same order. Strict comparison excludes a whole
timestamp tie; it never selects only part of one. Equal-size prefixes are
therefore identical. Under fixed $A$ and $B$ decisions, each retained final
event of the same degree has the same retained middle neighbors and the same
number of preceding retained $A$ events through each of those neighbors.
Its pair and triple contributions are identical.
\end{proof}

\begin{corollary}[Minimum-cost canonical final choices]\label{cor:canonical}
For each branch and each positive degree $j\le R$, order the optional $C$
events by nondecreasing cost, breaking cost ties by any fixed identifier order.
For each multiplicity $x_j$, it suffices to delete the cheapest $x_j$ identities
in that group. A feasible deletion never removes a $C$ event of degree $j>R$.
All optional degree-zero $C$ events may be retained when computing
minimum-cost count profiles.
\end{corollary}

\begin{proof}
The exchange in Lemma~\ref{lem:c} preserves both counts for every continuation.
Replacing a chosen size-$x_j$ subset by the cheapest such subset can only
decrease cost and never uses a forced-retained identity. The charge inequality
excludes degrees above $R$. A degree-zero final event participates in no
baseline pair and hence no triple; restoring it changes neither count and
cannot increase cost because deletion costs are nonnegative.
\end{proof}

If $m_{u,j}$ optional events of degree $j$ exist in branch $u$, enumerate
\begin{equation}
 \mathcal X_u=\left\{(x_j):0\le x_j\le m_{u,j},\quad
                    \sum_{j=1}^{R}jx_j\le R\right\}.
 \label{eq:configs}
\end{equation}
Only present positive-degree groups need explicit coordinates. Groups from
different branches are not merged. Configurations are processed one branch
at a time; no Cartesian product of all branches' choices is enumerated.

\begin{lemma}[Configuration bound]\label{lem:partitions}
Let
\[
 \Ppart=\bigl|\{(x_1,\ldots,x_R)\in\mathbb Z_{\ge0}^{R}:
                          \textstyle\sum_j jx_j\le R\}\bigr|,
 \qquad P_{\le0}=1.
\]
For $R>0$, $|\mathcal X_u|\le\Ppart\le
\exp(\pi\sqrt{2R/3})$. If $g_u$ positive-degree groups are present, a
depth-first enumeration visits at most $(g_u+1)\Ppart$ nodes.
\end{lemma}

\begin{proof}
The unrestricted multiplicity vectors are integer partitions of integers
between zero and $R$. Capacities remove vectors. For completeness, the
partition generating function gives, for every $\alpha>0$,
\[
 \Ppart\le e^{\alpha R}\prod_{j\ge1}(1-e^{-\alpha j})^{-1},\qquad
 \log\prod_{j\ge1}(1-e^{-\alpha j})^{-1}
 =\sum_{h\ge1}\frac{1}{h(e^{\alpha h}-1)}
 \le\frac{\pi^2}{6\alpha}.
\]
The last inequality uses $e^x-1\ge x$ and
$\sum_{h\ge1}h^{-2}=\pi^2/6$. Taking $\alpha=\pi/\sqrt{6R}$ proves the stated
bound, an elementary form of the classical partition estimate
\citep{hardy1918}. Enumerate each present group's count from zero to the
minimum of its capacity and the remaining loss allowance divided by its
degree. Every partial vector has a feasible zero extension. At a fixed
depth these extensions are distinct, so there are at most $\Ppart$ nodes
per depth. Counting the $g_u+1$ depths proves the traversal bound.
\end{proof}

For a fixed $x\in\mathcal X_u$, let $D_C(x)$ be the canonical final-event
deletions, $\Lambda_x=\sum_j jx_j$, and
$\kappa_x=\sum_{c\in D_C(x)}c_c$. For $b\in B_u$, set
\[
 s_b^x=|\{c\in C_u\setminus D_C(x):t(b)<t(c)\}|.
\]
The two-walk loss inside this branch is exactly
\begin{equation}
 \Lambda_x+\sum_{b\in D_0\cap B_u}s_b^x.
 \label{eq:disjointloss}
\end{equation}
This is disjoint accounting: first remove pairs ending at deleted $C$ events;
then remove pairs with deleted $B$ events and retained $C$ events. It remains
correct when both events of a baseline pair are deleted.


\subsection{Why unfinished configurations cannot be merged by scalar counts}

\begin{proposition}[Equal current statistics, different exact continuations]
\label{prop:unfinished}
There is a ten-event instance with two final-event deletion configurations
that have the same deletion cost, two-walk loss, and current three-walk count,
but for which retaining only either one as a representative prevents recovery
of the complete minimum-cost profile.
\end{proposition}

\begin{proof}
Use one root child $u$ and four children $v_1,\ldots,v_4$ of $u$.
Retain $B_i:v_i\to u$ at time $2i$, for $1\le i\le4$.
Make $C_j:u\to z$ at time $2j+1$ optional, for $1\le j\le4$.
There are two optional first events, $A_2:s_2\to v_2$ and
$A_4:s_4\to v_4$, both at time one. All costs are one. The baseline counts
are $(N_2,N_3)=(10,4)$. Set $T_0=5$, so $R=5$.

The final-event degrees are $1,2,3,4$. Consider
$D_{14}=\{C_1,C_4\}$ and $D_{23}=\{C_2,C_3\}$, before either first event
has been deleted. Both cost two, lose five two-walks, and leave two
three-walks. Their suffix-count vectors on $B_1,\ldots,B_4$ are respectively
$(2,2,1,0)$ and $(2,1,1,1)$. Completing the optional $A$ decisions gives:
\begin{center}
\begin{tabular}{lccc}
\toprule
Final-event context & Cost at $W_3=2$ & Cost at $W_3=1$ & Cost at $W_3=0$\\
\midrule
$D_{14}$ & 2 & $\infv$ & 3\\
$D_{23}$ & 2 & 3 & 4\\
Pointwise minimum & 2 & 3 & 3\\
\bottomrule
\end{tabular}
\end{center}
Costs include final-event deletions. Since all middle events must survive,
exact two-walk loss five can arise only from the degree subsets $1+4$ and
$2+3$. First-event deletions change no two-walk count. Keeping only $D_{14}$
loses the attainable coordinate $W_3=1$; keeping only $D_{23}$ gives an
incorrect minimum cost at $W_3=0$. These are the only two relevant final
contexts at that exact loss, proving the claim.
\end{proof}

This example rules out the stated one-representative compression. It is not
a lower bound on all possible state representations or algorithms. In the
algorithm below, each final-event configuration remains fixed until its child
tables are complete; pointwise minimization is then safe.

\section{An exact chronological dynamic program}\label{sec:dp}

Fix one branch configuration $x$ and a child $v$. Sort $A_v\cup B_v$ by time,
processing \emph{all $B$ events before all $A$ events at the same timestamp}.
Let $J_i(q,\ell,w)$ be the minimum optional deletion cost among decisions on
the first $i$ events such that:
\begin{itemize}
\item $q$ is the number of retained $A_v$ events processed so far;
\item $\ell$ is the sum of $s_b^x$ over processed deleted $B_v$ events;
\item $w$ is the number of triples contributed by processed retained $B_v$
events, using retained $C_u\setminus D_C(x)$ events.
\end{itemize}
An unattainable state has value $\infv$. Initialize only $J_0(0,0,0)=0$.
Every transition writes to a fresh layer. The four possible transitions are
listed below; a deletion is allowed only for an optional event.

\begin{center}
\begin{tabular}{llll}
\toprule
Event & Decision & New state & Added cost\\
\midrule
$a\in A_v$ & retain & $(q+1,\ell,w)$ & $0$\\
$a\in A_v$ & delete & $(q,\ell,w)$ & $c_a$\\
$b\in B_v$ & retain & $(q,\ell,w+q s_b^x)$ & $0$\\
$b\in B_v$ & delete & $(q,\ell+s_b^x,w)$ & $c_b$\\
\bottomrule
\end{tabular}
\end{center}

Reject a transition with $\ell>R-\Lambda_x$. All $w$ values remain distinct.
There is no rule that discards a larger $w$ because a smaller value was found
at the same or a lower cost. The final child table is
\[
 Q_{v,x}(\ell,w)=\min_q J_{|A_v|+|B_v|}(q,\ell,w).
\]
Each finite state stores one attaining original-ID deletion mask alongside
its cost. A keep transition preserves the mask; a deletion adds that event's
bit; a lower-cost replacement replaces both value and witness.

\begin{lemma}[Exactness of the child table]\label{lem:dp}
For fixed $x$, the recurrence computes the minimum optional $A_v\cup B_v$
deletion cost for every attainable exact pair $(\ell,w)$ with
$0\le\ell\le R-\Lambda_x$, with a valid witness.
\end{lemma}

\begin{proof}
Use induction on the chronological position. For an $A$ event, retention adds
one to the retained prefix count and deletion adds its cost. It affects no
root-ending two-walk. For a $B$ event, the tie rule ensures that exactly $q$
retained $A$ events have strictly earlier timestamps. Retaining it contributes
exactly $q s_b^x$ triples. Deleting it destroys precisely its $s_b^x$ remaining
pairs, as in~\eqref{eq:disjointloss}. Forced-retained events have only the keep
transition. These cases cover all legal decisions and each has the stated
effect.

At a fixed processing position, histories with the same $(q,\ell,w)$ have
identical possible future decisions and effects. Every previously processed
retained $A$ event precedes every future eligible $B$ event. Within a shared
timestamp, no $A$ event has been processed before its tied $B$ events, so this
claim also covers ties. Only the number of past retained $A$ events is needed
for a future prefix; past $B$ contributions have already been included in $w$.
Future event costs and forced statuses are fixed by the input and position.
Consequently, keeping only the minimum cost at an identical state is sound.
All loss increments are nonnegative, so a discarded loss above the cap can
never become feasible. The induction proves both attainability and minimum
cost. At the end, no future event uses $q$, which justifies minimizing over it.
Witness updates follow the same legal transitions.
\end{proof}

For two tables over disjoint event domains, define exact min-plus convolution
by
\begin{equation}
 (P\otimes Q)(\ell,w)=
 \min_{\substack{\ell_1+\ell_2=\ell\\w_1+w_2=w}}
        \{P(\ell_1,w_1)+Q(\ell_2,w_2)\}.
 \label{eq:convolution}
\end{equation}
Union the witnesses when a minimizing pair is chosen. For a fixed configuration,
convolve its child tables, then shift loss by $\Lambda_x$ and cost by
$\kappa_x$, and adjoin $D_C(x)$ to the witness. Take the pointwise minimum over
configurations \emph{after} each
configuration's complete branch table has been formed. Finally convolve
the complete branch tables. Loss above $R$ may be discarded throughout.

\begin{proof}[Correctness part of Theorem~\ref{thm:main}]
Every table transition, shift, and convolution constructs a legal deletion
set in disjoint original-event domains. Equations~\eqref{eq:identities}
and~\eqref{eq:disjointloss} show that its reported counts and cost are exact.
Thus every finite output entry has a valid attaining witness.

Conversely, take a minimum-cost feasible deletion set for any count pair.
Corollary~\ref{cor:canonical} replaces its final-event choices by canonical
choices without changing either count and without increasing cost. Each
branch's multiplicities occur in its enumerated family by
Lemma~\ref{lem:partitions}. Lemma~\ref{lem:dp} represents every remaining
child choice at no higher cost for its exact local counts. Children and
branches have disjoint event domains, and every root-ending pair or triple
belongs to one branch and one middle-event child. Therefore the corresponding
convolutions include the chosen exact global pair. The output cannot cost
more than its optimum, while its valid witness prevents it from costing less.
Adding $F$ and its fixed cost once completes the result.
\end{proof}

\section{Running time, bit complexity, and space}\label{sec:complexity}

Let $N=M+1$ and $S=(R+1)(N_3+1)$. Then $S=O(N^5)$ because $R\le M^2$
and $N_3\le M^3$. A child layer contains at most $NS$ states. Processing an
event uses a constant number of transitions per state. A table convolution
uses at most $S^2$ candidate pairs, even without an optimized convolution
algorithm. Minimizing over terminal $q$ uses at most $NS$ work per child.

There are at most $M$ events, children, and branches in their respective
disjoint collections. Each branch is processed for at most $\Ppart$
configurations. Consequently the complete recurrence uses
\begin{equation}
 O\!\left(\Ppart\,N[NS+S^2]\right)
 \label{eq:ops}
\end{equation}
table, arithmetic, and comparison operations, plus polynomial preprocessing
and configuration bookkeeping. The enumeration bound in
Lemma~\ref{lem:partitions} and the sum of the numbers of present groups keep
the same single factor $\Ppart$. Sorting timestamps and computing suffix
counts never enumerate the timestamp universe. Formula~\eqref{eq:ops} is at
most $O(\Ppart N^{11})$ operations; it is a conservative bound, not a practical
speed estimate.

Every finite cost is a sum of at most $M$ input costs and therefore has
$O(L)$ bits. Count indices have $O(\log(M+1))$ bits, and an original-ID witness
mask has $M$ bits. On a bit-cost random-access model, the arithmetic,
comparisons, and mask operations used here admit an $O(L+1)$ bound per
operation. Including polynomial preprocessing gives the deliberately loose
bound $O(\Ppart(L+1)^{12})$. Other conventional machine models may add a
polynomial factor without changing the $2^{O(\sqrt R)}\poly(L)$ conclusion.
There is no state dimension or loop up to the monetary budget or total cost.
Large costs affect the length of arithmetic, not the number of profile keys.

Process configurations sequentially and discard an individual configuration's
working tables after merging it into the branch aggregate. Retain only two
child layers, a current child table, current convolution tables, the aggregate
branch table, and the aggregate global table. A depth-first configuration
generator keeps polynomially many frames. Storing a witness mask in each
finite table entry therefore uses polynomial space. Output itself is also
polynomially bounded. This completes the complexity proof of
Theorem~\ref{thm:main}.

% Companion hardness section.
% Assumes amsmath, amssymb, amsthm, and the theorem/lemma/proof environments.
% No custom macros. Construction labels are prefixed hard:.
% The appended classical corollary uses the label full:np-complete.
% The complete-construction subsection can be moved to an appendix unchanged.

\section{Hardness parameterized by the deletion budget}
\label{hard:section}

The deletion budget and the preservation slack lead to different
parameterized complexity statements. The following companion result holds for the all-optional, unit-cost model. An original event is an
individually identified temporal contact; different contacts may have the
same endpoints and timestamp. Let \(E\) be the event set. For a retained
set \(S\subseteq E\), the quantity \(W_j(S)\) counts \(j\)-event walks
using events in \(S\), ending at the root, whose timestamps increase
strictly. Deleting \(D\) therefore leaves count \(W_j(E\setminus D)\).

\begin{theorem}[Hardness in the deletion budget]
\label{hard:budget}
Given an explicitly listed temporal in-arborescence of depth at most three,
deciding whether
\[
 |D|\le b,\qquad W_2(E\setminus D)\ge T_2,\qquad
 W_3(E\setminus D)\le T_3
\]
is \(\mathrm{W}[1]\)-hard parameterized by \(b\), even when every original
event is optional and has unit deletion cost. The restriction remains hard
with a single depth-one vertex and exactly one incoming depth-three event,
at timestamp zero, at each depth-two vertex that has such a predecessor.

The reduction from unweighted bipartite partial vertex cover with side
budgets \(a,q\) has \(b=a+2q\). Composed with the published reduction from
\(t\)-Clique, this gives \(b=t^2\).
\end{theorem}

\begin{proof}
The source is unweighted bipartite partial vertex cover with two separate
side budgets \cite[Definition~5 and Theorem~10]{mkrtchyan2024}.
The journal theorem states \(\mathrm{W}[1]\)-hardness for either side
budget; its displayed reduction from regular \(t\)-Clique uses
\(a=t\) and \(q=\binom{t}{2}\). Since both budgets are functions of
\(t\), that reduction also establishes hardness parameterized by
\(a+q\). The conference version is \cite[Theorem~6]{mkrtchyan2020}.
The regularization in the proof of Corollary~\ref{full:np-complete}
also preserves clique size, so the reduction may start from ordinary Clique.
The construction below has three steps.
Isolated rows and a pair of columns per source right vertex enforce the
two side budgets using the two temporal quotas.
Lemma~\ref{hard:retained-equivalence} proves this equivalence when the
middle events are temporarily required to survive.
Lemma~\ref{hard:replication} replaces that requirement by explicitly listed,
optional unit-cost events. Lemma~\ref{hard:size} proves polynomial encoded
size and the asserted parameter maps. These facts give a parameterized
many-one reduction and prove the theorem.
\end{proof}


The complete construction is in Appendix~\ref{hard:construction}.
Corollary~\ref{full:np-complete} gives NP-completeness of this decision
problem. It is a special case of the profile task: use unit costs, no
forced statuses, and select an entry with $w\le T_3$ and cost at most $b$.
The hardness parameter is $b$; no lower bound in preservation loss $R$ follows.


\section{Exact classifier optimization and its boundary}\label{sec:classifier}

\subsection{Predicates of the two counts}

\begin{corollary}[Count-based acceptance]\label{cor:classifier}
Let $g(W_2,W_3,z_0)$ be an acceptance predicate, where $z_0$ is proven invariant
under every allowed deletion, including forced deletions. After computing
$C$, minimum-cost acceptance is obtained by selecting the cheapest finite
entry with $g(N_2-\ell,w,z_0)=\mathrm{true}$. If one evaluation takes $t_g$
time, the additional bit time is $O(S(t_g+L+1))$, including cost comparisons.
In particular, a polynomial-time
predicate with polynomially encoded model data preserves the parameterized
bound, with those encodings included in the input length.
\end{corollary}

\begin{proof}
All legal deletion sets with a fixed count pair have the same acceptance
value. Replacing one by the minimum-cost witness at that pair preserves
acceptance. Completeness of the table then reduces the optimum to a scan
over its finite entries. A total monetary cap can be checked during the scan.
\end{proof}

For a classifier counterfactual, first compute the factual prediction on the
\emph{original history $E$}, before forced deletions. Set acceptance to the
desired class, which must differ from that factual class. A class change
already caused by mandatory deletions is legitimate and has the fixed forced
cost. The post-forced baseline is an algorithmic normalization, not a new
factual observation.

The predicate may be nonmonotone and may be implemented by a decision tree or
an ensemble. Its other features must actually be invariant. Replaying one
minimum-cost witness per pair through a model with additional dynamic features
can establish that witness's feasibility, but not full-model optimality:
another witness at the same pair can have different dynamic feature values.
A feature computed on the full history also cannot generally be obtained
by adding a constant to a windowed walk count, because a deleted event may
participate in a walk that crosses the window boundary. The present theorem
does not establish exactness for a classifier whose additional inputs change
under the allowed deletions.

\subsection{An elementary amount-feature obstruction}

Binary deletion costs are handled as minimized values. Binary-encoded dynamic
features are a different extension, because their attained values can define
the acceptance constraint.

\begin{proposition}[One additive dynamic feature can be hard at $R=0$]
\label{prop:amount}
If acceptance may additionally depend on the sum of retained binary-encoded
amounts on optional $A$ events, deciding whether a desired class is attainable
is NP-complete even when the collapsed static graph is a single depth-three
path, with $R=0$, unit deletion
costs, and a three-leaf classifier whose thresholds are part of the input.
\end{proposition}

\begin{proof}
Reduce from Subset Sum with positive integers $a_1,\ldots,a_n$ and target
$k$, restricted to $0<k\le S_a=\sum_i a_i$; trivial targets can be decided
before reduction. The standard source problem and its numerical input
dependence are discussed by Bringmann~\citep{bringmann2017}.

Use four vertices $r\to v\to u\to z$. Create $n$ optional parallel $A$ events
$r\to v$ at time zero, one forced-retained $B$ event $v\to u$ at time one,
and one forced-retained $C$ event $u\to z$ at time two. Give every optional
deletion cost one. Set $T=1$. Every legal choice has $W_2=1$, so $R=0$.
If an upper quota or budget is desired, $W_3\le n$ and budget $n$ are
nonbinding. Attach amount $a_i$ to optional event $i$ and define the dynamic
feature
\[
 H(D)=\sum_{i:\,A_i\notin D}a_i.
\]
The classifier accepts exactly when $H(D)=S_a-k$. A three-leaf tree implements
this on integer inputs using thresholds $S_a-k-\tfrac12$ and
$S_a-k+\tfrac12$, accepting only the middle interval. Its rational threshold
encodings have polynomial length. The original history has $H=S_a$ and is
rejected, since $k>0$.

A deletion set is accepted precisely when its deleted amounts sum to $k$.
The graph, events, amounts, and tree all have polynomial encoding length.
Membership in NP follows by checking a deletion witness, its walk counts,
the amount sum, and the tree output in polynomial time.
\end{proof}

This is an elementary Subset Sum boundary, not a claim of a new reduction
technique. The argument does not establish strong NP-hardness, does not give
a hardness result for the precise saved application model, and does not
contradict Theorem~\ref{thm:main}. It rules out an unrestricted $R$-only FPT
extension to this extra feature unless $\mathrm{P}=\mathrm{NP}$.

\section{Worked examples}\label{sec:examples}

All examples below use the original root-ending definition. Their standalone
JSON files are in the accompanying study artifact's
\texttt{implementation/examples} directory; the manuscript-source archive
contains the typesetting files.

\subsection{A nonmonotone target needs the complete profile}

On one path, put optional $A$ events at times $0$ and $2$, both of cost one;
force retention of $B$ events at times $1$ and $3$ and a $C$ event at time $4$.
Set $T=2$. The baseline has $(W_2,W_3)=(2,3)$ and $R=0$.

\begin{center}
\begin{tabular}{cccl}
\toprule
$W_2$ & $W_3$ & Minimum cost & Optional deletion witness\\
\midrule
2 & 3 & 0 & none\\
2 & 2 & 1 & later $A$ event at time 2\\
2 & 1 & 1 & earlier $A$ event at time 0\\
2 & 0 & 2 & both $A$ events\\
\bottomrule
\end{tabular}
\end{center}

Consider a hypothetical classifier whose desired class occurs only at
$W_3=2$, implemented by thresholds $1.5$ and $2.5$. The cheapest desired-class
change deletes the later $A$ event for cost one. An earliest-$A$ exchange
minimizes $W_3$ for a fixed number of unit-cost deletions, but selects the
other event at deletion count one. An algorithm that retains only that
minimum-$W_3$ value loses the accepting state. This example explains the
need for the exact $w$ coordinate. It is an illustrative classifier, not a
reported result on any trained model.

\subsection{Weighted decisions and mandatory deletions}

With optional $A$ events at times $0$ and $2$ of costs $100$ and $1$, and
retained $B,C$ events at times $3,4$, preserving $W_2=1$ gives $W_3$ values
$2,1,0$ at minimum costs $0,1,101$. The cheapest way to attain $W_3=1$
deletes the later event; a unit-cost earliest-event exchange is invalid here.

For the final-event exchange, retain $A$ at time $0$ and $B$ events at times
$1,3$, and make two $C$ events at time $4$ optional with costs $100,1$.
Both have degree two. With $T=2$, the attainable pairs are $(4,4)$ at cost
zero and $(2,2)$ at cost one. Canonicalization must select the cheap final
identity, even if an arbitrary identifier order lists the expensive one first.

Finally, retain $A$ at time $0$, force deletion of $B$ at time $1$ for cost
seven, and leave $C$ at time $2$ optional with cost three. Set $T=0$. The original counts
are $(1,1)$, the normalized baseline counts are $(0,0)$, and the minimum
profile has total cost seven with only the mandatory $B$ deletion. The
surviving $A$ event retains its original depth-three role. Deleting the now
degree-zero $C$ event cannot improve the minimum cost.

\section{Related work and parameter transfers}\label{sec:prior}
\subsection{Temporal reachability, counting, and cut problems}

Molter, Renken, and Zschoche~\citep[Theorem 14]{molter2024} give
polynomial-time algorithms for minimizing temporal reachability by deleting
or delaying time-edges when the underlying graph is a forest. Their
Lemma 2 transfers deletion to delay while preserving the forest property.
The stronger weighted forest delay formulation
\citep[Theorem 15, pp.~10--12]{molter2024} permits nonnegative weights on
reached vertices and undelayable time-edges, with a cardinality bound on
the modified time-edges. Its dynamic program records a vertex, a delay
budget, its first arrival time, and whether that arrival is supplied from
inside or outside its rooted subtree; a table entry minimizes the total
weight of reached vertices. The underlying temporal graph is undirected,
and the forest result allows arbitrary nonnegative traversal times,
including the strict convention with traversal time one. Thus tree
structure, strictness, and protected contacts alone are not novelty
distinctions, and the source's vertex weights must not be confused with
arbitrary event-deletion costs.

The final journal version also gives a $4^x|G|^{O(1)}$ algorithm for both
reachability problems in the non-strict setting, where $x$ is the timed
feedback vertex number~\citep[Theorem 16]{molter2024}. This parameter is
zero on underlying forests. The output contract here instead requires
minimum original-event deletion cost at every attainable exact pair of
root-ending two- and three-event tuple counts, subject to the preservation
quota and shared event decisions. A derivation from the reachability
results must therefore supply an encoding or sufficient state for these
multiplicities, costs, and constraints, as well as the claimed dependence
on preservation loss. Appendix~\ref{app:reachability} gives a direct-summary
counterexample: the same reached sets and earliest arrivals can coexist
with different tuple counts. It does not exclude richer states or auxiliary
reductions. The journal article extends the earlier conference version
\citep{molter2021}. The comparison above uses the final journal
statements and proofs, with their journal theorem numbering.

Enright, Meeks, and Molter~\citep[Theorem 2]{enright2025counting} count temporal
paths on forests in $O(|V|\tau^2)$ time, where $\tau$ is the lifetime. Their
Theorems 6--7 treat further structural parameters. Fixed-deletion counting
does transfer to the present setting: Appendix~\ref{app:counting} gives a single
normalization preserving strict ties, event identities, and both counts
under every deletion subset, with lifetime at most $M$. Thus strictness,
orientation, and parallel contacts do not obstruct that counting transfer.
The additional task here is to choose one shared deletion set and compute
minimum costs for the entire attainable profile with the stated dependence
on preservation loss.

Robust Temporal Cut~\citep{enright2026robust} removes contacts so that no
source--target temporal path survives even after bounded timestamp
perturbations. For undirected temporal graphs, its stated tractability
parameters include timed vertex
cover, vertex cover plus solution size, and temporal neighborhood diversity
plus solution size. This is a related temporal-deletion problem with a
different feasibility condition. The existence of those algorithms does not
itself supply an exact reduction for two simultaneous multiplicity counts.
Any proposed transfer must also preserve its controlling parameters.

\subsection{Cover formulations and the parameters they retain}

There is a direct exact set-system representation. Each baseline two-walk is
a red element, each baseline three-walk is a blue element, and each optional
event supplies the set of elements containing its original identifier.
Selecting event sets incurs their deletion costs. A profile entry
$C(\ell,w)$ requires red union coverage exactly $\ell$ and blue union
coverage exactly $N_3-w$. An upper quota on surviving triples instead
requires a partial blue target. Shared incidences matter: a walk destroyed
by several deletions is still counted once.

This representation relates the problem to red--blue set cover. Abidha and
Ashok~\citep[Theorem 6 in the author preprint]{abidha2024} give an exact
fixed-dimensional hyperplane algorithm parameterized by red coverage for
covering all blue points. The temporal witness incidence admits a
three-dimensional coordinate-plane representation, but the compared
feasible family has no independent charge for selected planes and freely
selects red-free sets. First-event choices in the temporal model can be red-free and still
cost money. The partial-blue approximation setting of
Chlamt\'a\v{c}, Makarychev, and Vakilian~\citep[Appendix A]{chlamtac2023}
is another close precedent for the benefit--damage formulation; its
approximation guarantee does not by itself decide exact profile feasibility
and minimum event cost.

Madireddy, Nandy, and Pandit~\citep[Theorems 1--2]{madireddy2025} give exact
algorithms for intervals and for a restricted axis-parallel-line case. The
direct temporal witness representation need not be an interval system.
With one first event $a$, two middle events $b_1,b_2$, and two final events
$c_1,c_2$, all temporally compatible, the four triples $q_{ij}=(a,b_i,c_j)$
induce the event sets
\[
 \{q_{11},q_{12}\},\quad \{q_{21},q_{22}\},\quad
 \{q_{11},q_{21}\},\quad \{q_{12},q_{22}\}.
\]
An interval order would need four distinct adjacent pairs on four blue
points, which have only three adjacencies. The restricted line algorithm
also relies on disjoint red incidences, whereas a temporal red pair belongs
to both its middle and final events. These observations concern the direct
representation and do not rule out another exact encoding.

The bipartite partial-cover framework of
Mkrtchyan and coauthors~\citep{mkrtchyan2024} includes a selected-vertex cost
budget and covered-edge profit targets. Their hardness theorem with two side budgets
provides the starting problem for the reduction in Appendix~A. For the positive comparisons,
neither selected cost nor covered profit is inherently bounded by the preservation loss $R$.
For example, take $n$ optional unit-cost first events before one retained
middle/final pair and require preservation of that pair. Then $R=0$ while
destroying all triples requires $n$ deletions. A target leaving at most
$\lfloor n/2\rfloor$ triples makes both the destroyed-triple target and the
surviving-triple allowance grow at $R=0$. An algorithm controlled by those
quantities therefore needs an additional parameter argument to yield the bound in Theorem~\ref{thm:main}. This is a statement about the indicated transfer, not about all
possible reductions.

\subsection{Generalized deletion propagation}\label{sec:gdp}

Generalized deletion propagation (GDP) permits simultaneous deletion and
preservation requirements on multiple query outputs, with shared decisions
on source tuples and an objective expressed through view-deletion counts
\citep[Definition 4.1]{makhija2025gdp}. Its ILP framework supports self-joins
and both set and bag semantics; its tractability results concern specified
deletion-propagation variants and query conditions
\citep[Sections 5--6]{makhija2025gdp}.

Appendix~\ref{app:gdp} derives a polynomial-size GDP encoding of each
unit-cost temporal profile entry. Protected compatibility relations compile
strict timing and root incidence; retaining every event identifier in the
query output preserves tuple multiplicity under set semantics. Deletion and
preservation bounds on the same query enforce an exact count. Mandatory
statuses are handled by normalization and singleton preservation queries.
Polynomially many target pairs recover the full profile table. This is a
derived encoding, not a temporal theorem stated in the cited work.

The appendix also gives a weighted witness-indicator ILP with the original
binary cost coefficients. That objective extends the literal cardinality
objective of GDP Definition~4.1; expanding an event into one copy per cost
unit would not justify polynomial input size. Thus the formulation and
full-table output are not the claimed novelty. The contribution claimed in
Theorem~\ref{thm:main} is the structural algorithm with its joint time and
space guarantee. Obtaining that guarantee through GDP requires a solving-time
and parameter-transfer argument beyond the encoding.


\subsection{General representations and optimization frameworks}
Courcelle and Durand's fly-automata framework computes satisfying assignments,
cardinality spectra and related functions (Theorem 23), and develops minimizing
algebras and weighted cardinality optimization in Section 3.2.3
\citep{courcelle2016}. It is a direct precedent for retaining sufficient
attainable information and optimizing within states. Its complexity discussion
requires bounds on the chosen automata and representations. A transfer needs
an explicit representation preserving shared event decisions and exact
counters, together with bounds implying the stated running time. The quotient and enumeration in Sections~3--4 provide one such construction. A suitable application to
this class could provide another presentation of the algorithm.

Del Pia proves polynomial solvability for factorized binary polynomial
optimization with fixed structural parameters, with explicit dependence on
the factorization in Theorems 1--2 and a conversion result in Proposition 2
\citep{delpia2026}. This framework can express the weighted profile search.
The issue is the parameter transfer, not expressibility. The following direct
encoding makes that distinction concrete.

For an exact target pair $(p,q)$ and a trial optional monetary cap $b$, use
one Boolean deletion variable per optional event. Substitute forced statuses
and express $W_2(x)$ and $W_3(x)$ as sums of retained-pair and retained-triple
monomials in $(1-x_e)$. Let $C_O=\sum_{e\in O}c_e$, take $0\le b\le C_O$, and
encode a nonnegative slack $\sigma$ using enough binary digits to represent
every integer through $C_O$ (no digits are needed if $C_O=0$). Maximize
\begin{equation}
 -\left[(W_2(x)-p)^2+(W_3(x)-q)^2+
      \left(b-\sum_{e\in O}c_ex_e-\sigma\right)^2\right].
 \label{eq:penalty}
\end{equation}
Its maximum is zero exactly when the pair is attainable within cap $b$.
After Boolean multilinearization, the degree is at most six and the expanded
representation has polynomial size and coefficient bit length. First test
$b=C_O$ to detect an unattainable coordinate. For an attainable one, binary
search over $b$ requires only $O(1+\log(C_O+1))$ calls, including the zero-cost
case. If a polynomial format omits constants, carry the known offset into
the comparison threshold; constant or purely linear cases can be solved
directly. The encoding above is a derivation for the temporal problem;
the cited paper does not state this transfer. A fixed degree
does not by itself give a bounded number of factors. The cited runtime's
factor-dependent input exponents therefore do not, without another bound,
yield $2^{O(\sqrt R)}\poly(L)$ for this transfer.

Kundu and coauthors' Theorem 1.7 concerns constrained, class-based MaxSAT with a true-variable
budget and a randomized approximation guarantee \citep{kundu2026}. Unit-cost
quota encodings can introduce event-deletion variables and pair-loss variables;
their budget includes both the deletion allowance and $R$. For integral
class totals, sufficiently small approximation error makes a successful
returned solution meet the original targets exactly, retaining the theorem's
success probability of at least $1-1/e$. The resulting dependencies need not
be bounded by $R$: the deletion allowance and objective target can grow already
at $R=0$. The weighted exact-profile extension also needs an explicit monetary
constraint encoding. The unit-cost budget argument cannot simply be relabeled
as arbitrary event-cost support.

Khouna and coauthors formulate optimal tree-ensemble counterfactual search
using a desired prediction, a cost, and plausibility/actionability constraints
(Section 2, Equation 1), then compare exact solution paradigms
\citep{khouna2026}. Thus targeting an actual class and minimizing a change cost
are not new objectives. The classifier corollary applies the specialized temporal-profile algorithm
only when the model depends on the supported counts and invariant inputs.
It does not introduce desired-class optimization or a general framework for
counterfactual search.


\subsection{Scope of the comparison}

The inspected results supply important counting, representation, and
optimization tools. The GDP construction in Appendix~\ref{app:gdp}
establishes a semantic transfer for the unit-cost formulation and gives an
explicit weighted ILP extension. It does not analyze the time needed to
solve those encoded instances as a function of preservation loss.
An algorithmic transfer yielding Theorem~\ref{thm:main} must preserve the
shared original-event decisions and exact output contract, and its
representation size and parameter dependence must together imply
$2^{O(\sqrt R)}\poly(L)$ time and polynomial space, including the complete
input encoding and an absolute polynomial exponent.

The direct proof supplies this guarantee for the stated temporal class.
The comparison does not prove that no alternative derivation or earlier
result can imply the theorem. Exact profiles, weighted minima, integer
partitions, min-plus convolution, and terminal predicate filtering are not
claimed as new general techniques.

\section{Implementation and validation}\label{sec:validation}
\subsection{Comparison with exhaustive enumeration}

For each small test instance, an exhaustive checker enumerates ordered pairs
and triples of original events, requiring matching endpoints, strictly
increasing timestamps, and a final endpoint at the root. For every legal
deletion set, it counts the surviving tuples without calling the solver's
counting, grouping, or optimization routines. At each quota it compares
the attainable count pairs and their minimum costs with those returned by
the solver, and checks every deletion witness directly. Tied optimal
witnesses need not use the same event identifiers.

The frozen corpus uses seed 20261104. It contains 20 named histories,
all $3^5=243$ timestamp assignments from $\{0,1,2\}$ on a fixed five-event
path, 80 seeded random histories, and 32 transformed histories. The
transformations are input permutation, injective identifier renaming,
order-preserving affine timestamp transformation, and positive cost scaling
on eight selected histories. Named cases include zero costs, strict ties,
parallel events, forced disconnection, multiple branches, and unequal costs
within equal-degree final-event classes.

\begin{center}
\begin{tabular}{lr}
\toprule
Check & Completed\\
\midrule
Frozen histories & 375\\
Quota queries / valid-input solver calls & 1,863\\
Feasible / infeasible quota queries & 1,062 / 801\\
Queries with $R=0$ & 561\\
Legal deletion subsets enumerated & 81,108\\
Exact profile entries and witnesses checked across queries & 2,529\\
Acceptance/budget-selection queries & 54,076\\
Returned selection witnesses checked & 19,731\\
Invalid inputs correctly rejected & 57\\
Transformation relations checked & 32\\
\bottomrule
\end{tabular}
\end{center}

Selection tests use nine predicates, including intermediate exact counts,
parity and disconnected count regions, and monetary caps at the optimum and
adjacent values. Witness verification checks original identities, forced
statuses, optional and total costs, both root-ending counts, and the quota.
The largest history has 12 events and costs of up to 203 bits. Maximum
baseline counts are $W_2=21$ and $W_3=6$. The profile total counts checks
across quota queries, not necessarily distinct global coordinates.

\subsection{Compatibility and native-run evidence}

The original full suite passed on CPython 3.12.14 on Linux. A subsequent
compatibility repair removed Python-version-sensitive dataclass and evaluated
type-alias syntax. All 14 solver function abstract syntax trees were unchanged.
Smoke checks covered CPython 3.8.20 and 3.9.25 on Linux, and the full frozen
suite passed on CPython 3.9.25 with identical reference output digests.

The macOS validation summary records a successful run of the same suite on
4 November 2026 using CPython 3.9.6 on macOS 27.0.1 (ARM64).
The SHA-256 hashes of the solver, validator, corpus, and protocol match
those of the compatibility-corrected release. The output hashes recorded
in the summary match those computed from the retained reference files.
For this suite, the archive contains the macOS summary but not the
full macOS outputs; those files could not be hashed separately.

The implementation requires Python 3.8 or newer and the standard library.
These are finite correctness and compatibility checks. No runtime-performance
comparison is inferred from their elapsed times. The general result rests
on the proof. The implementation archive also retains the initial root-endpoint
specification correction, both frozen oracle versions, and their audit
records; Appendix~\ref{app:reproducibility} explains the provenance.

The ten-event illustration in Proposition~\ref{prop:unfinished} is archived
separately with its stated expected profiles and a direct enumeration of all
64 legal optional-deletion subsets. It is an explanatory check added with
this manuscript and is not counted in the original 375-history suite.

\subsection{Exact targets in synthetic transaction histories}\label{sec:batch}

The file \texttt{HI-Small\_Trans.csv} from version~8 of IBM's synthetic
anti-money-laundering dataset~\citep{altman2023} contains 5,078,345 transactions.
The sampling unit was a destination
account and transaction minute. An earlier audit fixed 1,024 queries by
seeded hash order from the 3,854,758 distinct units with a full 24-hour
lookback, without using labels or structural eligibility. These queries
contain 1,023 distinct roots. Each context includes every transaction in
$[t-1440,t)$ between accounts within incoming static distance three of the
root in that window. The extraction retained all such events, including loops, tied
times and parallel contacts; no conflicting edge was removed to make a
context eligible. The cohort had been studied previously, and queries may
share events. The results describe these queries rather than independent
trials or an unseen test set.

\paragraph{Coverage.}
All 1,024 contexts were extracted and their baseline counts determined.
The structural check was applied to the original induced context.
Table~\ref{tab:batch-coverage} retains empty, unsupported, and
resource-limited queries in the denominator. Of 232 queries with
baseline $W_3>0$, 77 fit the model and 68 received complete profiles.
The nine skipped eligible positive queries had a median of 91 events
compared with 12 for the 68 profiled queries.

\begin{center}
\begin{tabular}{lrrr}
\toprule
Query disposition & $W_3=0$ & $W_3>0$ & Total\\
\midrule
Empty context & 179 & 0 & 179\\
Outside the required tree structure & 125 & 155 & 280\\
Eligible, above the 32-event cap & 34 & 9 & 43\\
Eligible, all requested profiles completed & 454 & 68 & 522\\
\midrule
Total & 792 & 232 & 1,024\\
\bottomrule
\end{tabular}
\refstepcounter{table}\label{tab:batch-coverage}\\[3pt]
{\small Table \thetable. Disposition of all sampled account/minute queries.}
\end{center}

\paragraph{Two questions asked of the same profile.}
Every event was optional with unit deletion cost. Given baseline counts
$N_2,N_3$, set $T=\lceil pN_2/100\rceil$ for $p\in\{100,95,90\}$ and
$U=\lfloor 3N_3/4\rfloor$. For each $T$, the evaluation compared the goals $W_3\le U$ and
$W_3=U$, both subject to $W_2\ge T$. Budgets 0, 1 and 3 were applied to
the minimum deletion cost. The 454 profiled queries with $N_3=0$ satisfy
both goals without intervention; the comparison below therefore concerns
the 68 profiled queries with $N_3>0$.

\begin{center}
\begin{tabular}{lrrr}
\toprule
 & \multicolumn{2}{c}{Attainable within three deletions}
 & Unattainable at any cost\\
\cmidrule(lr){2-3}\cmidrule(l){4-4}
Preservation & $W_3\le U$ & $W_3=U$ & $W_3=U$\\
\midrule
$100\%$ & 68 & 13 & 55\\
$95\%$  & 68 & 14 & 54\\
$90\%$  & 68 & 15 & 53\\
\bottomrule
\end{tabular}
\refstepcounter{table}\label{tab:batch-goals}\\[3pt]
{\small Table \thetable. Each row refers to the same 68 positive, fully profiled queries.}
\end{center}

At full preservation, the least cost for the upper bound was one deletion in
59 queries, two in eight, and three in one. In 55 of the 68, no deletion set
could attain the exact same target while preserving every pair, even without
a deletion budget. At $R=0$, every middle or final event belonging to a
baseline pair must survive. Deleting a first event then removes all triples
extending it as one group; different first events remove disjoint groups.
Attainable triple losses are sums of these group sizes, so some intermediate
integer losses can be absent. For example, in the preidentified pilot
query \texttt{q0011}, preserving seven pairs permits either two triples at
cost zero or zero triples at cost one. Exactly one triple is unattainable.
Rounding also matters: on 50 of the 68 positive queries all three
percentages require the same integer $T$. The rows of
Table~\ref{tab:batch-goals} are paired comparisons, not 204 independent
observations.

\paragraph{Verification and scope.}
Deduplicating equal values of $T$ gave 653 planned tables among the 565
eligible nonempty queries. The 32-event cap excluded 90 tables before
optimization; these have unknown outcomes. There were no solver or oracle
timeouts. The Mac run reported witness replay for all 935 returned entries
across the 563 completed tables. For 528 of those tables (at most 16
optional events), exhaustive enumeration agreed on every attainable pair
and its minimum cost. The other 35 had witness replay but no exhaustive
comparison. The run provides evidence about local count sensitivity and
model coverage on synthetic transactions. It does not measure classification
accuracy, changes in fraud decisions, or comparative runtime.


\section{Discussion and conclusion}\label{sec:conclusion}

Preservation loss bounds the structural choices of deleted final events even
when the deletion budget, cost magnitudes, and number of branches are not
parameter-bounded. The continuation equivalence of final events makes those
choices enumerable by integer partitions. Conditional chronological states
then recover exact minimum-cost profiles, and branchwise aggregation avoids
multiplying configuration families. The full proof yields an absolute
polynomial exponent and polynomial space, including original-event witnesses.

The GDP encoding clarifies the scope of this result. Simultaneous exact
count requirements and a complete table of count optima fit a broader
optimization framework; the structural state, charging argument, and
preservation-loss complexity bound are the contribution proved here.
The weighted ILP formulation provides another exact representation of the
same feasible event choices. Neither that encoding nor the fractional
example in Appendix~\ref{app:gdp} rules out stronger formulations or
another derivation through existing methods.

The budget-hardness result concerns a different parameter and establishes no
tight lower bound for the square-root dependence on $R$. The algorithm's
polynomial factor is deliberately conservative; the finite checker is a
correctness artifact rather than an efficiency experiment. A practical
comparison with a generic exact ILP baseline should use the same event
identities, costs, forced statuses, and full profile output, with controlled
variation in $R$. Such a runtime comparison is not part of the present
evaluation.

The supported model has explicit events, fixed depth at most three, strict
timestamps, nonnegative deletion costs, and root-ending walk counts. Negative
costs, compressed event multiplicities, deeper temporal structures, or extra
dynamic feature constraints require separate analysis. An independent
cardinality cap in addition to arbitrary costs is not encoded in the current
table. These boundaries leave precise questions for further work while keeping
the present theorem and its classifier corollary fully specified.

\appendix
% BEGIN COMPLETE CONSTRUCTION: this block may be placed in an appendix.
\section{Complete construction for budget hardness}
\label{hard:construction}

\paragraph{Source problem.}
An instance of two-budget bipartite partial vertex cover consists of a
simple bipartite graph \(H=(X,Y,E_H)\), nonnegative integer budgets \(a,q\),
and an integer target \(\ell\). It asks for sets \(X'\subseteq X\) and
\(Y'\subseteq Y\), with \(|X'|\le a\) and \(|Y'|\le q\), such that at least
\(\ell\) edges have an endpoint in \(X'\cup Y'\).
Write \(n_0=|X|\), \(m=|Y|\), and \(e=|E_H|\).
The source is unweighted; the two budgets are separate cardinality caps.

\paragraph{The regular-Clique source map.}
\label{hard:clique-source}
The following paragraphs detail the preceding reduction in the cited
article and its boundary cases. Let \(G\) be a simple \(d\)-regular graph with \(N\)
vertices, \(m_G\) edges, and nonnegative Clique target \(t\). Cases
\(t<3\), \(t>N\), and \(t=N\) are decidable in polynomial time; the
last case tests completeness. Cases \(d<t-1\) or
\(m_G<\binom{t}{2}\) are no-instances. Map every directly decided case
to the three-event temporal yes/no template below, carrying the numerical
budget \(b=t^2\). It remains to consider
\(3\le t<N\), \(d\ge t-1\), and \(m_G\ge\binom{t}{2}\).

Make one vertex of \(X\) for each vertex of \(G\), and one vertex of
\(Y\) for each edge of \(G\). Join \(v\in X\) to \(g\in Y\) exactly
when \(v\) is not an endpoint of \(g\). This non-incidence graph has
\(N+m_G\) vertices and \(e=(N-2)m_G\) edges. Set
\[
 a=t,\qquad q=\binom{t}{2},\qquad
 \ell=t(m_G-d)+(N-t)q.
\]
Any source selection within these caps can be extended to exactly
\(t\) left vertices and \(q\) right vertices without losing covered
edges. For these saturated choices \(S\subseteq V(G)\) and
\(J\subseteq E(G)\), let
\(i(S,J)=\sum_{g\in J}|\operatorname{endpoints}(g)\cap S|\).
Inclusion--exclusion gives the exact covered-edge count
\begin{equation}
\label{hard:clique-coverage}
\begin{aligned}
 C(S,J)
 &=t(m_G-d)+q(N-2)-\bigl(tq-i(S,J)\bigr)\\
 &=\ell-\bigl(2q-i(S,J)\bigr).
\end{aligned}
\end{equation}
Since \(i(S,J)\le2q\), reaching \(\ell\) is equivalent to every edge
of \(J\) having both endpoints in \(S\). A simple graph on \(t\)
vertices has at most \(q=\binom{t}{2}\) edges, so every such \(S\) is a
clique. Conversely, a \(t\)-clique and its internal
edges attain equality. This is the reduction of
\cite[Theorem~10]{mkrtchyan2024}, now with its at-most-to-exact step and
coverage identity explicit. The construction is polynomial, both side
budgets depend only on \(t\), and the temporal map below gives
\(b=a+2q=t^2\). If the resulting partial-cover target exceeds \(e\),
the following preprocessing correctly returns a no-instance.

\paragraph{Trivial partial-cover instances.}
The cases \(\ell\le0\) and \(\ell>e\) are immediate. If \(a\ge n_0\) or
\(q\ge m\), selecting an entire side covers every edge. If \(q=0\), choose
the \(a\) largest left degrees, or all left vertices when fewer are
available: edges incident to different left vertices are distinct, so this
maximizes coverage. Thus the construction may assume
\begin{equation}
\label{hard:source-range}
 0\le a\le n_0,\qquad 1\le q<m,\qquad 0\le\ell\le e.
\end{equation}
To preserve the exact formula \(b=a+2q\) on the directly decided cases as
well, use the three-event path
\(s\mathbin{\to}v\mathbin{\to}u\mathbin{\to}z\), at times \(0,1,2\),
with the same numerical budget \(b=a+2q\). Set \((T_2,T_3)=(1,1)\)
for a yes-instance and \((T_2,T_3)=(2,1)\) for a no-instance.
The first accepts \(D=\varnothing\); the second is impossible because its
undeleted two-walk count is one. Every event in either instance is
optional.

\paragraph{Rows, numerical constants, and columns.}
Add \(h=2m+1\) isolated left vertices to \(H\), and number the resulting
rows \(1,\ldots,n\). Define
\begin{equation}
\label{hard:constants}
\begin{aligned}
 n&=n_0+2m+1, & r&=n-a, & b&=a+2q,\\
 \kappa&=e+1, & P&=4m^2+1,\\
 \Lambda&=\bigl((\kappa+1)+r\kappa\bigr)mP+e+1 .
\end{aligned}
\end{equation}
In particular,
\begin{equation}
\label{hard:gaps}
\begin{gathered}
 r>2m,\qquad \kappa>e,\qquad P>2m^2,\\
 \Lambda+1>(\kappa+1)mP,\qquad
 \Lambda>r\kappa mP+e.
\end{gathered}
\end{equation}
The isolated rows have no source edges and require no additional budget.

For each \(y_j\in Y\), introduce an original column of numerical weight
\(j\) and a companion column of numerical weight \(P-j\).
All \(2m\) weights are distinct: the original weights are \(1,\ldots,m\),
and \(P-m>m\). Sort the columns by increasing weight and write their weights
as \(\rho_1<\cdots<\rho_{2m}\). Their sum is \(mP\).
These weights specify event counts in the construction; they are not event
deletion costs.

For row \(i\) and column \(t\), let \(N_{it}\) be the adjacency indicator
of the padded graph when column \(t\) is original. Set \(N_{it}=0\) for
companion columns and for padded rows. Let
\(\sigma_t=\sum_{i=1}^n N_{it}\); hence \(0\le\sigma_t\le n_0\).

\paragraph{The intermediate temporal instance.}
Use the static arcs
\[
 s_i\longrightarrow v_i\longrightarrow u\longrightarrow z
 \quad (1\le i\le n),\qquad v_0\longrightarrow u.
\]
The root is \(z\). There is a single depth-one vertex \(u\); the auxiliary
depth-two vertex \(v_0\) has no predecessor. For each row \(i\), put one
event \(A_i:s_i\to v_i\) at time zero. For every sorted column \(t\), put one
event \(C_t:u\to z\) at time \(2t\).

Prescribe the following cumulative numbers of middle events strictly
before \(C_t\):
\begin{equation}
\label{hard:prefixes}
 \mu_{it}=\Lambda+\kappa\rho_t+N_{it}
 \quad(1\le i\le n),\qquad
 \mu_{0t}=n+n\rho_t-\sigma_t .
\end{equation}
Set \(\mu_{i0}=0\) for \(0\le i\le n\).
For each \(t\), put exactly \(\mu_{it}-\mu_{i,t-1}\) distinct events
on \(v_i\to u\), all at time \(2t-1\), for every \(0\le i\le n\).
Call these the \(B\)-events.

All these multiplicities are nonnegative integers. The initial prefixes
are positive. For consecutive columns after the first, their increments
are bounded below by
\[
 \kappa(\rho_t-\rho_{t-1})+(N_{it}-N_{i,t-1})
 \ge\kappa-1\ge0
\]
on a productive row, and by
\[
 n(\rho_t-\rho_{t-1})-(\sigma_t-\sigma_{t-1})
 \ge n-n_0>0
\]
on the auxiliary row. The times \(2t-1\) and \(2t\) realize precisely the
prefixes in \eqref{hard:prefixes}, with strict temporal inequalities.
All contacts, including contacts sharing endpoints and a timestamp, receive
distinct IDs.

For the moment, \(B\)-events must survive, while \(A\)- and \(C\)-events are
optional. This is an intermediate construction only; the final instance
will make every event optional.

\begin{lemma}[Exact counts in the intermediate instance]
\label{hard:counts}
Let \(I\subseteq\{1,\ldots,n\}\) be the retained \(A\)-rows and
\(J\subseteq\{1,\ldots,2m\}\) the retained \(C\)-columns. Put
\[
 p=|I|,\qquad s=|J|,\qquad
 \omega=\sum_{t\in J}\rho_t,\qquad
 N(I,J)=\sum_{\substack{i\in I\\t\in J}}N_{it}.
\]
When all \(B\)-events survive, the counts and deletion cost are
\begin{equation}
\label{hard:identities}
\begin{aligned}
 W_2^\circ(I,J)
   &=n\bigl[s(\Lambda+1)+(\kappa+1)\omega\bigr],\\
 W_3^\circ(I,J)
   &=p\bigl[s\Lambda+\kappa\omega\bigr]+N(I,J),\\
 |D_{AC}|&=(n-p)+(2m-s).
\end{aligned}
\end{equation}
\end{lemma}

\begin{proof}
Each two-walk ending at \(z\) consists of a \(B\)-event followed by a
retained \(C_t\). Its contribution at column \(t\) is
\[
 \sum_{i=1}^n\mu_{it}+\mu_{0t}
 =n\Lambda+n\kappa\rho_t+\sigma_t
    +n+n\rho_t-\sigma_t
 =n\bigl[(\Lambda+1)+(\kappa+1)\rho_t\bigr].
\]
An \(A\)-event cannot belong to a two-walk ending at the root.
Each three-walk consists of \(A_i,B,C_t\), with \(i\in I\), \(t\in J\),
and the \(B\)-event counted by \(\mu_{it}\). The unique \(A_i\) precedes
every \(B\)-event on its row. The auxiliary row contributes no three-walks.
Summing \(\mu_{it}\) over \(I\times J\) proves the second identity. The last
identity counts the deleted optional IDs.
\end{proof}

Define the target retained-column count and target weight by
\[
 s_\star=2(m-q),\qquad \omega_\star=(m-q)P.
\]
They satisfy \(s_\star\ge2\), \(\omega_\star>0\), and
\(n+2m-b=r+s_\star\). Set the intermediate quotas to
\begin{equation}
\label{hard:intermediate-quotas}
\begin{aligned}
 T^\circ
   &=n\bigl[s_\star(\Lambda+1)+(\kappa+1)\omega_\star\bigr],\\
 U^\circ
   &=r\bigl[s_\star\Lambda+\kappa\omega_\star\bigr]+e-\ell.
\end{aligned}
\end{equation}

\begin{lemma}[The two quotas enforce the budget allocation]
\label{hard:allocation}
An \(A/C\) deletion set of cost at most \(b\) satisfies
\(W_2^\circ\ge T^\circ\) and \(W_3^\circ\le U^\circ\) if and only if
\begin{equation}
\label{hard:forced-values}
 p=r,\qquad s=s_\star,\qquad \omega=\omega_\star,\qquad
 N(I,J)\le e-\ell.
\end{equation}
Every feasible intermediate deletion set has cost exactly \(b\).
\end{lemma}

\begin{proof}
Suppose the two quotas hold. Since \(0\le\omega\le mP\), if
\(s<s_\star\) then
\[
 \frac{W_2^\circ-T^\circ}{n}
 =(s-s_\star)(\Lambda+1)
     +(\kappa+1)(\omega-\omega_\star)
 \le-(\Lambda+1)+(\kappa+1)mP<0,
\]
contradicting the lower quota. Therefore \(s\ge s_\star\).

Suppose instead that \(s=s_\star+d\), where \(d\ge1\).
The deletion budget gives \(p+s\ge r+s_\star\), so \(p\ge r-d\).
Since \(d\le s\le2m<r\), the latter lower bound is positive. Moreover,
\[
 ps\ge(r-d)(s_\star+d)
    =rs_\star+d(r-s)\ge rs_\star+1.
\]
All terms in the formula for \(W_3^\circ\) are nonnegative. It follows that
\[
\begin{aligned}
 W_3^\circ
 &\ge (rs_\star+1)\Lambda\\
 &>rs_\star\Lambda+r\kappa mP+e\\
 &\ge rs_\star\Lambda+r\kappa\omega_\star+e-\ell
   =U^\circ,
\end{aligned}
\]
a contradiction. Consequently \(s=s_\star\).

The budget now gives \(p\ge r\), and the lower quota gives
\(\omega\ge\omega_\star\). Subtracting the base term of \(U^\circ\)
from the triple count yields
\[
\begin{aligned}
 W_3^\circ-r(s_\star\Lambda+\kappa\omega_\star)
  ={}&(p-r)(s_\star\Lambda+\kappa\omega_\star)\\
    &+p\kappa(\omega-\omega_\star)+N(I,J).
\end{aligned}
\]
All three terms on the right are nonnegative, and their sum is at most
\(e-\ell\le e\). If \(p>r\), the first term is at least
\(s_\star\Lambda\ge2\Lambda>e\). If \(\omega>\omega_\star\), integrality
makes the second term at least \(r\kappa>e\).
Thus \(p=r\), \(\omega=\omega_\star\), and the remaining inequality is
\(N(I,J)\le e-\ell\). Conversely, substituting
\eqref{hard:forced-values} into \eqref{hard:identities} verifies both quotas.
The same substitution gives cost \((n-r)+(2m-s_\star)=a+2q=b\).
\end{proof}

\begin{lemma}[Companion balance]
\label{hard:balance}
If \(s=s_\star\) and \(\omega=\omega_\star\), exactly \(m-q\) original
columns and \(m-q\) companion columns survive.
\end{lemma}

\begin{proof}
Let \(c\) be the number of retained companion columns. Their weights and
the original-column weights give
\[
 \omega=cP+\delta,\qquad
 \delta=\sum_{\substack{j:\ \text{original }j\\\text{survives}}}j
       -\sum_{\substack{j:\ \text{companion }j\\\text{survives}}}j.
\]
There are at most \(2m\) summands, each of absolute value at most \(m\),
so \(|\delta|\le2m^2<P\). As \(\omega=(m-q)P\), the integer
\((m-q-c)P=\delta\) can only be zero. Hence \(c=m-q\);
the remaining \(s_\star-c=m-q\) retained columns are original.
The argument does not require matching identities for retained originals
and companions.
\end{proof}

\begin{lemma}[Equivalence with retained middle events]
\label{hard:retained-equivalence}
The source partial-cover instance is feasible if and only if the
intermediate temporal instance admits a deletion set of cost at most
\(b\), with \(W_2^\circ\ge T^\circ\) and \(W_3^\circ\le U^\circ\).
\end{lemma}

\begin{proof}
First take a source cover \(X',Y'\). Extend its selected left vertices to
exactly \(a\) rows of the padded graph, and extend \(Y'\) to exactly \(q\)
right vertices. Both extensions are possible and preserve coverage.
Delete the corresponding \(a\) events \(A_i\). For every selected right
vertex \(y_j\), delete its original \(C\)-event and its matching companion
\(C\)-event. This uses exactly \(a+2q\) deletions.
The retained columns consist of \(m-q\) pairs, each of weight \(P\).
Thus \(p=r\), \(s=s_\star\), and \(\omega=\omega_\star\).
The quantity \(N(I,J)\) is exactly the number of source edges not covered
by the extended selections, so it is at most \(e-\ell\).
Lemma~\ref{hard:allocation} gives the two temporal quotas.

Conversely, consider a feasible intermediate deletion set.
Lemma~\ref{hard:allocation} gives exactly \(a\) deleted \(A\)-events, and
Lemma~\ref{hard:balance} gives exactly \(q\) deleted original \(C\)-events
and \(q\) deleted companion \(C\)-events. Select the original left vertices
whose \(A\)-events were deleted, omitting any deleted padded rows, and
select the right vertices whose original \(C\)-events were deleted.
These selections use at most \(a\) left vertices and exactly \(q\) right
vertices. Only original columns and original rows can have nonzero
\(N_{it}\). Therefore \(N(I,J)\) counts precisely the source edges with
neither endpoint selected. Since \(N(I,J)\le e-\ell\), the selections
cover at least \(\ell\) edges.
\end{proof}

\paragraph{Removing every retention requirement.}
Set
\begin{equation}
\label{hard:replication-parameters}
 Q=2mb+1,\qquad T_2=QT^\circ,\qquad T_3=QU^\circ.
\end{equation}
Replace each intermediate \(B\)-event by \(Q\) distinct optional events
with the same endpoints and timestamp. Keep the \(A\)- and \(C\)-events
unchanged, and retain budget \(b\). All final events are original input
IDs with unit deletion cost. The intermediate instance is used only to
prove this final construction; it imposes no retention restriction on the
final input.

\begin{lemma}[Optional replication]
\label{hard:replication}
The all-optional final instance is feasible if and only if the
intermediate instance is feasible. In fact, every feasible final deletion
set deletes no \(B\)-event.
\end{lemma}

\begin{proof}
Fix the \(A/C\) choices of an arbitrary final deletion set \(D\) with
\(|D|\le b\), and let \(w_j\) be the intermediate count with those
choices and all middle events retained. Restoring every deleted \(B\)
copy in the final instance gives exactly \(Qw_j\) surviving \(j\)-walks,
for \(j=2,3\).

A \(B\) copy belongs to at most \(2m\) surviving three-walks: its row has
at most one \(A\)-event, and there are at most \(2m\) final \(C\)-events.
An auxiliary copy belongs to none. Each triple contains exactly one
\(B\) copy, so the triple loss caused by all deleted \(B\) copies is at
most
\[
 2m\,|D\cap B|\le2mb=Q-1.
\]
If \(w_3>U^\circ\), integrality would therefore imply
\[
 W_3(E\setminus D)\ge Q(U^\circ+1)-(Q-1)=QU^\circ+1>T_3,
\]
contrary to feasibility. Hence \(w_3\le U^\circ\).
Restoring middle events can only increase the two-walk count; consequently
\[
 Qw_2\ge W_2(E\setminus D)\ge QT^\circ,
\]
and \(w_2\ge T^\circ\).
The \(A/C\) part of \(D\) alone costs at most \(b\), so it is a feasible
intermediate solution. By Lemma~\ref{hard:allocation}, that part already
uses exactly \(b\) deletions. Thus \(D\) contains no deleted \(B\) copy.

Conversely, a feasible intermediate \(A/C\) deletion set can be used in
the final instance while deleting no \(B\) copy. The two counts are then
exactly \(Qw_2\) and \(Qw_3\), and both final quotas hold.
\end{proof}

\begin{lemma}[Encoded size and parameter map]
\label{hard:size}
The nontrivial construction contains exactly
\begin{equation}
\label{hard:event-count}
 M=n+2m+Qn\bigl[\Lambda+(\kappa+1)P-\kappa\bigr]
\end{equation}
explicitly listed events, uses timestamps in \(\{0,\ldots,4m\}\), and
has size polynomial in the explicit source input length.
Its deletion budget is \(b=a+2q\), and the composition with the
\(t\)-Clique reduction in \cite[Theorem~10]{mkrtchyan2024} has \(b=t^2\).
\end{lemma}

\begin{proof}
The last column has weight \(P-1\) and is a companion; hence all its
adjacency indicators and its \(\sigma\)-value are zero. Each productive
row contains \(\Lambda+\kappa(P-1)\) intermediate \(B\)-events, and the
auxiliary row contains \(nP\). Their total is
\[
 n\bigl[\Lambda+\kappa(P-1)\bigr]+nP
   =n\bigl[\Lambda+(\kappa+1)P-\kappa\bigr].
\]
Multiplying by \(Q\) and adding \(n\) events of type \(A\) and \(2m\)
of type \(C\) gives \eqref{hard:event-count}.
The static tree has \(2n+3\) vertices and depth three. Its only
depth-one vertex is \(u\). Replication changes neither the static tree
nor any timestamp.

Put \(S_H=n_0+m\). In the range \eqref{hard:source-range},
\[
\begin{gathered}
 e=O(S_H^2),\quad n,r,b=O(S_H),\quad
 \kappa,P,Q=O(S_H^2),\\
 \Lambda=O(S_H^6),\quad M=O(S_H^9).
\end{gathered}
\]
Moreover \(T^\circ,U^\circ=O(S_H^8)\), so \(T_2,T_3=O(S_H^{10})\).
Event IDs, endpoints, timestamps, budgets, and quotas therefore need
only \(O(\log S_H)\) bits per numerical field. In particular, an explicit
event list has length \(O(S_H^9\log S_H)\). Every contact is written out;
the construction uses no binary-encoded multiplicities.
The directly decided cases have constant-size temporal event lists and
retain the binary encoding of \(b=a+2q\).

The parameter bound \(b=a+2q\le2(a+q)\) is independent of \(n_0,m,e\).
The source theorem uses budgets
\(a=t\) and \(q=\binom{t}{2}\) in its reduction from regular
\(t\)-Clique. Thus the composed budget is
\[
 b=t+2\binom{t}{2}=t^2.
\]
This identity also holds if a directly decided source case uses the
three-event template described above.
\end{proof}

\paragraph{Why the preservation-slack algorithm is consistent with this result.}
For the nontrivial construction, the final instance has no mandatory
deletions and \(T_2>0\).
Its preservation slack \(R=W_2(E)-T_2\) is exactly
\begin{equation}
\label{hard:slack}
 R=Qnq\bigl[2(\Lambda+1)+(\kappa+1)P\bigr].
\end{equation}
Indeed, the full retained column set has \(s=2m\) and weight \(mP\);
subtracting \(QT^\circ\) from the first identity in
\eqref{hard:identities}, scaled by \(Q\), gives \eqref{hard:slack}.
This quantity is unbounded in the source graph size even with fixed
budgets. For example, fix \(a=q=1\), \(m=2\), \(e=\ell=1\), and increase
the number of isolated left vertices. Then \(b=3\) is fixed, whereas
\(n\), \(\Lambda\), and \(R\) grow.
Thus this reduction does not bound \(R\) by a function of \(b\).
It is consistent with a \(2^{O(\sqrt R)}\operatorname{poly}(|\mathcal I|)\)
algorithm parameterized by preservation slack. The conclusion is
\(\mathrm{W}[1]\)-hardness in \(b\); no \(\mathrm{W}[1]\)-membership claim
or ETH exponent lower bound is asserted.
% END COMPLETE CONSTRUCTION.

\subsection{Classical complexity consequence}

The explicit construction also yields the following classical consequence.
The proof is included for completeness, without a historical-priority claim
for this corollary or the elementary regularization used below.

\begin{corollary}[Classical NP-completeness]
\label{full:np-complete}
The decision problem in Theorem~\ref{hard:budget} is NP-complete, even with
all events optional and of unit cost, one depth-one vertex, and exactly
one time-zero \(A\) event at each productive depth-two vertex.
\end{corollary}

\begin{proof}
Membership in NP follows from an original-event deletion bit vector.
Its budget and forced statuses, when present, can be checked directly.
Enumerating ordered retained event pairs and triples, and testing their
endpoints and strict timestamp inequalities, computes the required counts
using \(O(M^3)\) candidate tuples. Endpoint and timestamp comparisons,
comparison with the encoded quotas, and validation of the structural
restrictions all take polynomial time in the complete input length.

Classical hardness follows by a polynomial chain of reductions from ordinary Clique.
Let \(G\) have \(N\) vertices and maximum degree \(\Delta\), with target
\(t\). Cases \(t<3\), \(t>N\), or \(\Delta=0\) are directly decidable.
Otherwise make \(2\Delta\) disjoint copies of \(G\), named \(L_i,R_i\)
for \(i\in\mathbb Z_\Delta\). For every original vertex \(v\) of degree
\(d(v)\), join its copy in \(L_i\) to its copies in \(R_{i+j}\), for
\(0\le j<\Delta-d(v)\), with indices modulo \(\Delta\).
Each copied vertex receives exactly its degree deficit, on both sides,
so the resulting simple graph is \(\Delta\)-regular. No added edge lies
in a triangle: its endpoints have the same original label in different
copies; a vertex of another label can be an internal neighbor of at most
one endpoint, while all edges among copies of the same label are
bipartite. Thus every clique of size at least three remains within an
original copy, preserving the existence of a \(t\)-clique.
The regularization has \(2N\Delta=O(N^2)\) vertices and
\(N\Delta^2=O(N^3)\) edges and is polynomial.

The regularized graph has \(N'=2N\Delta>t\), so the strict source-size
condition is automatic. Apply the regular-Clique prechecks and source-map
argument above. For its \(N'\) vertices, \(m'\) edges, and degree \(d\),
the non-incidence construction in \cite[Theorem~10]{mkrtchyan2024}
uses one source-side vertex per vertex and one per edge, with
\[
 a=t,\qquad q=\binom{t}{2},\qquad
 \ell=(N'-t)\binom{t}{2}+t(m'-d).
\]
The graph has \(N'+m'\) vertices and at most \(N'm'\) edges, so this
step is polynomial. Apply the temporal construction above.
Its side-size sum satisfies \(S_H=O(N^3)\); hence
Lemma~\ref{hard:size} gives polynomial encoded output size, even with
the loose event bound \(O(N^{27})\). Directly decided source cases use
the stated constant-event yes/no templates. The output has the promised
structural restrictions, and its budget is
\(b=t+2\binom{t}{2}=t^2\).
This polynomial many-one reduction proves NP-hardness and completes
the proof.
\end{proof}

Consequently, an exact polynomial-time algorithm for every instance of
this full model would imply \(\mathrm P=\mathrm{NP}\). This consequence
does not rule out a simpler or faster algorithm parameterized by
preservation slack \(R\), and it supplies no matching lower bound for
the square-root exponent.

\section{Counting-model transfer and a reachability-summary distinction}

\subsection{A simultaneous normalization for every deletion subset}
\label{app:counting}

The forest counting theorem of Enright, Meeks, and
Molter~\citep{enright2025counting} uses undirected temporal paths with
nondecreasing timestamps. The following explicit transformation explains
why fixed-deletion counting transfers to the present directed, strict,
parallel-event convention. The transformation below establishes the comparison;
it is not asserted as a theorem stated in the cited work.

\begin{proposition}[Deletion-preserving counting normalization]
For an explicit depth-three temporal in-arborescence with $M>0$ events, there
is a polynomial-time transformation to a temporal forest with at most $M$
distinct time labels such that, simultaneously for every retained event
subset, the number of temporal paths from vertices of original depth $h$
to the root is exactly $W_h$, for $h=2,3$. Original event identities are
preserved bijectively.
\end{proposition}

\begin{proof}
Let $\rho(e)$ be the rank of $t(e)$ among distinct original timestamps,
preserving ties. Let $d(e)\in\{1,2,3\}$ be the original depth of the source
of $e$, and give each original event a distinct index
$i(e)\in\{0,\ldots,M-1\}$. Define
\[
 t'(e)=4M\rho(e)+Md(e)+i(e).
\]
For consecutive rootward events $e,f$, their depths differ by one, so
\[
 t'(e)-t'(f)=4M(\rho(e)-\rho(f))+M+i(e)-i(f).
\]
The final term lies between $1$ and $2M-1$. Hence
$t(e)<t(f)$ if and only if $t'(e)\le t'(f)$. Equal original timestamps
are incompatible in the rootward direction. All transformed labels are
distinct: different ranks are separated by the $4M$ term, different depths
within a rank by $M$, and different identities within a rank and depth by
their indices. Rank-compressing these labels preserves their order and
leaves at most $M$ time labels.

Forget arc directions. The underlying graph is a tree, so the unique simple
path from any vertex to the root follows the original rootward direction.
A temporal path cannot leave that route and return without repeating a
vertex. Its successive transformed labels are nondecreasing exactly when
the original labels are strictly increasing. Distinct original parallel
contacts remain distinct time-edges. The correspondence survives deletion
of any common subset of identities and preserves the number of events in
each path. Summing over starting vertices of original depth $h$ gives the
claimed equality.
\end{proof}

Empty histories are handled directly. Forced statuses and costs may be
carried as identity metadata, but the cited counting theorem does not
optimize them. Counting for each fixed deletion set therefore transfers;
enumerating and optimizing a shared deletion set still requires a separate
argument. This normalization concerns simple temporal paths on a forest,
not arbitrary undirected walks that may return along an edge.

\subsection{Reached sets and earliest arrivals do not encode tuple quotas}
\label{app:reachability}

Consider the path $a\to v\to u\to z$, with three optional first events at
times $1,4,7$, three retained middle events at times $2,5,8$, and a retained
final event at time $10$. There are three two-walks and six three-walks.
Deleting the first event at time one or the first event at time four leaves,
respectively, three or four triples; both choices leave all three pairs and
have the same unit deletion cost.

Reverse every arc and replace time $t$ by $11-t$, taking $z$ as the source
and assigning traversal time one. For either deletion, the same four vertices
are reachable with earliest-arrival times $(0,2,4,5)$ along the reversed
path. The final reversed contact at time four comes from the undeleted
original first event at time seven in both cases. Yet the original upper
quota $W_3\le3$ accepts only the first deletion. No statistic consisting
solely of these reached sets and earliest arrivals can distinguish the two
choices, even after assigning fixed vertex weights.

This example separates the indicated summary from tuple multiplicities.
It does not prove that reachability methods cannot be extended, nor that an
auxiliary encoding is impossible.


\section{Reference implementation and local use}

\begin{minipage}{\linewidth}
The implementation exports two functions:
\begin{lstlisting}
solve_profiles(instance)
select_profile(result, accept, max_cost=None)
\end{lstlisting}
\end{minipage}

Its chronological state is exactly $(q,\ell,w)$; values are cost/mask pairs.
Its branch aggregation uses exact min-plus convolution. Sparse dictionaries
avoid materializing impossible states; no dominance rule across different
count coordinates is used. Cost ties may choose any valid minimum-cost witness.

The JSON result distinguishes \texttt{original} and \texttt{baseline} counts,
reports normalized \texttt{min\_w2} and \texttt{R}, and lists profiles sorted by
\texttt{(loss,w3)}. Each entry contains \texttt{w2}, \texttt{w3},
\texttt{optional\_cost}, \texttt{total\_cost}, and sorted
\texttt{deleted\_ids}. Invalid schemas or original graphs raise
\texttt{ValueError}; a quota above the baseline returns an infeasible result.
Unknown extra JSON fields are ignored; they do not become new features or
constraints. The CLI permits arbitrary-length integer decimal encodings by
temporarily disabling the interpreter's decimal-conversion limit when that
setting exists. Importing the module does not change interpreter settings.

The implementation uses the Python standard library and supports Python
3.8 or newer. From the repository's \texttt{temporal-profiles/implementation}
directory, run:
\begin{lstlisting}
python3 demo.py
python3 profile_solver.py examples/nonmonotone.json --output result.json
\end{lstlisting}
The implementation archive documents how to run the exhaustive checker and
reports the results. These examples run locally with the Python standard library.

\section{Reproducibility and evidence provenance}\label{app:reproducibility}

\subsection{Canonical artifacts}

The repository preserves the 5 November 2026 manuscript snapshot together
with the compatibility-corrected implementation, examples, frozen corpus and
validator, retained output files and summaries, the macOS validation summary,
and the separate ten-event illustration check. Its frozen manifest identifies
those artifacts. The present 12 November revision is distributed separately
with its own source and checksums; it adds the deletion-propagation comparison
while retaining the same algorithm, proofs of the existing results, and
reported experimental evidence. The original supplement and historical audit
records identify earlier sources and must not be confused with this
manuscript revision.

The corrected implementation and frozen input identities are:
\begin{lstlisting}
Solver:
a9f8f2c9f276055ab570940d784244192d177c16299f7a295ad75c481de7f4ba
Validator:
29ff8cfbe8df6f3aaedc9f7f0ae54c5a69a7081c805576a80d8f84760d5c8891
Corpus:
9a6979e08652fb949fa8eb7ca9d9a432204142c24e954d6fa046c172b631e67c
Protocol:
bcca092b7e763b8781fbdd261d3fff91f90c0dcdb70a965898f878de025e7297
\end{lstlisting}

The output digests recorded in the Mac summary and matched against the
retained reference bytes are:
\begin{lstlisting}
Candidate outputs:
faf9fc7f8960d46d81509ccbb711e83eb6c5a6657361b061113f7980131cad4a
Selection outputs:
133df11b935269866ea2766743a2b8037f0b88fd8b5be14bb2bf7743062e9f02
\end{lstlisting}

The macOS validation summary is identified by the following SHA-256 digest:
\begin{lstlisting}
26eb54eaebed908e15fe45371c48f7b9703eb8838b30ad9900c67af82d7f574f
\end{lstlisting}
The Linux compatibility checks and the macOS validation run are documented
separately in the implementation archive.

\subsection{Aggregate transaction-cohort evidence}

The separate Mac batch in Section~\ref{sec:batch} used the same solver hash
as the frozen correctness suite above. Its aggregate results archive
(\texttt{BATCH\_RESULTS.zip}) has SHA-256
\begin{lstlisting}
cd16912a10fe1f1e8444824f5a0648b57f2c2b05133e50be5bb1f24d522204c2
\end{lstlisting}
It contains the batch protocol, one result per query, one result per policy
condition, and machine- and human-readable summaries. Each query has 18
policy conditions (three preservation levels, two target comparisons, and
three budgets); repeated rounded lower quotas use the same table. The
reported source, cohort, protocol and program hashes match the frozen
configuration. The repository's frozen cohort file supplies the synthetic
bank/account identifiers and cutoff times needed to select the same queries.
The aggregate archive does not contain the original dataset, complete local
contexts, full profiles, or native deletion-witness diagnostics; those detailed
run files remained on the machine that ran the batch. The witness replays
and exhaustive comparisons reported here were performed on that machine.
Reaggregating the export checks outcome totals without independently
replaying those witnesses. This batch and the worked application pilot are
separate from the 375-history correctness suite.

\subsection{The retained endpoint-specification correction}

The initial checker specification omitted the requirement that a counted walk
end at the root. For a three-event path, the checker therefore counted both
length-two subwalks, although only one ends at the root. The solver already
enforced this requirement and did not change. The corrected checker tests
whether the final event ends at the root.

All 375 original histories and 1,779 original quota queries were preserved.
The amendment added 84 missing thresholds derived from the corrected root
baseline, yielding 1,863 queries. The archive retains the initial failure,
both corpus/protocol/oracle snapshots, exact query additions, and the amended
successful run. This correction changed the expected counts used for validation.

\subsection{Build and run separation}

The solver and the exhaustive checker run locally with the Python standard
library. The manuscript builds with \texttt{latexmk -pdf main.tex} from this
revision's \texttt{source} directory; the accompanying Mac build script places
generated files in a separate build directory. Rebuilding the PDF does not
rerun scientific tests. The implementation and validation archive and the Mac
cohort runner are distinct artifacts from this manuscript-source revision.

\paragraph{Research workflow.}
Goutam Dyapa ran the reported macOS experiments. ChatGPT assisted with
manuscript drafting and revision, code development, and validation tasks.

\section{Deletion-propagation encoding and a weighted formulation}
\label{app:gdp}

This appendix derives a representation of the temporal problem in the GDP
framework~\citep[Definition 4.1]{makhija2025gdp}, and a weighted ILP extension.
The construction concerns expressibility. It does not transfer the
running-time guarantee of Theorem~\ref{thm:main} to a general-purpose solver.

Use the normalization of Section~\ref{sec:model}. In particular, validate the
original temporal graph before applying forced deletions, set $E_0=E\setminus
F$, and reserve the fixed cost $c(F)$. If $T_0>N_2$, return infeasibility.
Otherwise let $R=N_2-T_0$. The original event identifiers are distinct even
when endpoints and timestamps coincide.

\subsection{Protected compatibility data and unit costs}

Create a unary relation $\mathsf{Event}(i)$ with one row for every identifier
in $E_0$. For $h\in\{2,3\}$, let $H_h$ contain exactly the ordered
$h$-tuples of identifiers forming strictly time-increasing walks ending at
the designated root in $E_0$. These relations have at most $M^2$ and $M^3$
rows. They can be constructed by polynomially many endpoint and
binary-timestamp comparisons; the numeric lifetime is not expanded.

Protect each nonempty $H_h$ by placing its identity query in the GDP
preservation list with threshold $|H_h|$. Since source deletion cannot add
rows to an identity view, no row of that relation can be deleted. An empty
relation needs no protection requirement. The remaining view queries are
ordinary conjunctive queries:
\begin{align*}
 Q_2(i,j)&\mathbin{:\!-}H_2(i,j),\mathsf{Event}(i),\mathsf{Event}(j),\\
 Q_3(i,j,k)&\mathbin{:\!-}H_3(i,j,k),\mathsf{Event}(i),
                 \mathsf{Event}(j),\mathsf{Event}(k).
\end{align*}
Strict comparisons and the root requirement are thus compiled into protected
relation membership, without assuming built-in inequality syntax in GDP.

\begin{proposition}[Unit-cost profile encoding]\label{prop:gdp-encoding}
Suppose all optional events have unit deletion cost. For every
$0\le\ell\le R$ and $0\le w\le N_3$, a GDP instance of size polynomial in
$L$ has a feasible solution exactly when $C(\ell,w)<\infv$; its optimum is
$C(\ell,w)$. An optimal source-deletion set yields an attaining original-event
witness in polynomial time.
\end{proposition}

\begin{proof}
Write $w_2=N_2-\ell$ and $w_3=w$. For each $h\in\{2,3\}$, put $Q_h$ in
the deletion list with requirement $N_h-w_h$ whenever this value is positive,
and in the preservation list with requirement $w_h$ whenever that value is
positive. The former requires at most $w_h$ surviving output rows; the latter
requires at least $w_h$. Omit every zero requirement, consistent with the
positive thresholds in GDP Definition~4.1. At the endpoints, the omitted
inequality is automatic because $0\le |Q_h|\le N_h$.

For each $e\in K$, preserve the Boolean singleton query
$\mathsf{Keep}_e()\mathbin{:\!-}\mathsf{Event}(e)$ with requirement one.
Here $e$ is the constant original identifier. Minimize deletions in the
identity view $Q_E(i)\mathbin{:\!-}\mathsf{Event}(i)$, with no maximization
views. The protected compatibility rows cannot be deleted, and these
singleton requirements prohibit deletion of events in $K$. Therefore every
feasible source-deletion set is exactly a set $D_0\subseteq O$ of optional
event rows.

For any such $D_0$, a row of $Q_h$ survives precisely when every event in its
identifier tuple belongs to $E_0\setminus D_0$. Its cardinality is therefore
$W_h(E_0\setminus D_0)$. All tuple identifiers are output coordinates, so set
semantics does not merge parallel-event walks. A single deleted event affects
all rows containing it in both views. Hence the two exact-count requirements
hold exactly for the deletion sets defining $C(\ell,w)$, and the objective
equals $|D_0|$. There are $O((M+1)^3)$ data rows and $O(M+1)$ view
requirements, with polynomially encoded identifiers and thresholds. This
proves the size, feasibility, optimum, and witness assertions.
\end{proof}

The full deletion witness is $F\cup D_0$ and its total cost is
$c(F)+C(\ell,w)$. The GDP objective above concerns optional deletions only;
forced cost remains the reserved constant.

\subsection{Binary costs and full-profile recovery}

For arbitrary nonnegative binary-encoded costs, use a binary retention
variable $x_e$ for every $e\in E_0$, fixed to one for $e\in K$. Introduce a
binary variable $z_q$ for each tuple $q\in H_2\cup H_3$, and impose
\begin{align}
 z_q&\le x_e &&(q\in H_2\cup H_3,\ e\in q),\label{eq:gdp-and-upper}\\
 z_q&\ge \sum_{e\in q}x_e-|q|+1
              &&(q\in H_2\cup H_3),\label{eq:gdp-and-lower}\\
 \sum_{q\in H_2}z_q&=N_2-\ell,\qquad
 \sum_{q\in H_3}z_q=w.\label{eq:gdp-exact}
\end{align}
The objective is
\begin{equation}
 \min\left(c(F)+\sum_{e\in E_0}c_e(1-x_e)\right).
 \label{eq:gdp-cost}
\end{equation}
For binary variables,~\eqref{eq:gdp-and-upper}--\eqref{eq:gdp-and-lower}
make $z_q$ the conjunction of its event-retention decisions. Thus this ILP
has optimum $c(F)+C(\ell,w)$ whenever the target is attainable, and is
otherwise infeasible. A walk destroyed by several deletions still contributes
only one lost tuple. There are $O((M+1)^3)$ variables and constraints, and
all coefficients have polynomial bit length in $L$. In particular, retaining
$c_e$ as a binary coefficient avoids replication proportional to its value.
This is an explicit weighted witness formulation, not a claim that arbitrary
binary event costs are the literal objective of GDP Definition~4.1.

Solving each target independently and retaining an optimal witness recovers
the full table. The number of target pairs is at most
\[
 (R+1)(N_3+1)\le (M^2+1)(M^3+1)=O((M+1)^5).
\]
Consequently, full-profile output is expressible through polynomially many
exact optimization calls. Their representation size does not bound their
solving time in terms of $R$. The particular tractability results in
\citet[Section 6]{makhija2025gdp} do not automatically transfer under these
joint exact-count requirements and cost coefficients; a proposed transfer
must verify the relevant hypotheses and parameter bounds.

\paragraph{An elementary relaxation diagnostic.}
Let $b_1,b_2:v\to u$ be distinct events at time zero, and let $c:u\to z$ be
an event at time one. All three have unit cost and are optional. This
depth-two instance has $(N_2,N_3)=(2,0)$. With $T_0=1$ and target
$(W_2,W_3)=(1,0)$, the integral optimum is one: delete either middle event;
deleting $c$ misses the target. In the relaxation of the displayed ILP to
$[0,1]$, the assignment
\[
 x_{b_1}=x_{b_2}=1,\qquad x_c=\tfrac12,\qquad
 z_{(b_1,c)}=z_{(b_2,c)}=\tfrac12
\]
is feasible and costs $\tfrac12$. Thus this relaxation is not automatically
integral on the temporal class. This diagnostic concerns the displayed
formulation; it excludes neither stronger formulations nor other algorithms.
It does not replace the structural proof of Theorem~\ref{thm:main} or prove
that no earlier result can imply its guarantee.


\bibliographystyle{plainnat}
\bibliography{references}
\end{document}
